\scanpage{4}
% L01–02
\deleted{Cas général \uline{OPS} $k$ noeth., et que $J$ idéal de carré nul, $k_0=k/J$, \unsure{$a_0=b_0$} $u_0={}$}

% L03–04
\uline{Lemme} Soit $a\in\deleted{\unsure{E}}k^n$, $u\in\operatorname{Aut}(G_{a\,k})$ tel que $u\circ f_a$ soit de la forme $f_b$ ($b\in k^n$), alors $u=\mathrm{id}_{G_{a\,k}}$.

% L05–06
\deleted{Cas général OPS $k$ noeth., $J$ idéal de carré nul et $u$ tel que $u_0=\mathrm{id}_{G_{a\,k_0}}$ ($k_0=k/J$). On a}
% L07
\[
u(x)=c_0x+c_1x^p+c_2x^{p^2}\ \text{---}+c_Nx^{p^N},\quad \unsure{c_i|a_i}\in J
\]
% L08–12
\[
\begin{aligned}
u\circ f_a(x)&\mathrel{{:}}c_0(x^{p^n}+a_{n-1}x^{p^{n-1}}\unsure{+|-}a_{n-2}x^{p^{n-2}}\ \text{---}+a_1x^p+a_0x)\\
&+c_1(x^{p^{n+1}}+a_{n-1}^{p}x^{p^n}\ \text{--}\qquad\text{--}+a_0^p x^p)\\
&+c_2(x^{p^{n+2}}+a_{n-1}^{p^2}x^{p^{n+1}}+\qquad +a_0^{p^2}x^{p^2})\\
&\text{---}\\
&+c_N({\unsure{x|\alpha}}^{p^{n+N}}+a_{n-1}^{p^N}x^{p^{n+N-1}}+\text{---}\qquad \unsure{-|+}a_0^{p^N}x^{p^N})
\end{aligned}
\]
% L13–15
Le coeff\textsuperscript{t} de $x^{p^{n+N}}$ dans l'expression totale est $c_N x^{p^{n+N}}$, donc \deleted{\unsure{pour}} si $N\geq 1$, on doit avoir $c_N=0$. Donc de proche en proche $c_{N-1}=0\ldots c_1=0$.

% L16–17
Donc $u(x)=c_0(x)=c_0 x^{p^n}\ \text{---}$, et il faut avoir $c_0=1$ cqfd.

% L18–19
c) $\varphi_n$ \deleted{\unsure{est}} \uline{propre} Critère valuatif, il faut prouver \deleted{\uline{\unsure{Corollaire|Conclusion}} \unsure{$X_{n\,\mathrm{réd}}$ est intègre}}

% L20–24
\uline{Lemme} Soit $S$ \inserted{$=\operatorname{Spec}(V)$} trait sur $\mathbb F_p$, \inserted{pt gén. $y$} $H\subset G_{a\,S}$, \inserted{$H\in X_n(S)$ i.e. $H$ plt fini $/S$ de rg $p^n$,} supposons que $H_y=\operatorname{Ker} f_{a_y}$, $\unsure{a_y|a|a_\eta}\in K^n$ ($K=k(y)$), alors $a_y\in V^n$ et \deleted{\unsure{$H=$} \unsure{$\operatorname{Ker}(f_a:G_{a\,S}$}} $H=\operatorname{Ker}(f_a:G_{a\,S}\mathbin{\text{---}}G_{a\,S})$.
% archive sheet number: 3
